% --- CORRECCIÓN IMPORTANTE EN LA PRIMERA LÍNEA ---
% Agregamos "xcolor=table" para que funcione \rowcolor sin errores
\documentclass[aspectratio=169, 10pt, xcolor=table]{beamer}

% --- Configuración del Tema (estilo profesional como el ejemplo) ---
\usetheme{Madrid}
\usecolortheme{dolphin}

% --- Paquetes necesarios ---
\usepackage[utf8]{inputenc}
\usepackage[spanish]{babel}
\usepackage{amsmath, amssymb}
\usepackage{booktabs}
\usepackage{graphicx}
\usepackage{listings}
\usepackage{tikz}
\usetikzlibrary{arrows.meta, positioning, shapes.geometric}

% --- Ruta de Imágenes (crea la carpeta "imagenes" y coloca ahí tus archivos) ---
\graphicspath{{./imagenes/}}

% --- Estilo para código Python ---
\definecolor{codegreen}{rgb}{0,0.6,0}
\definecolor{codegray}{rgb}{0.5,0.5,0.5}
\definecolor{codepurple}{rgb}{0.58,0,0.82}
\definecolor{backcolour}{rgb}{0.95,0.95,0.92}
\lstset{
    language=Python,
    backgroundcolor=\color{backcolour},
    commentstyle=\color{codegreen},
    keywordstyle=\color{blue},
    numberstyle=\tiny\color{codegray},
    stringstyle=\color{codepurple},
    basicstyle=\ttfamily\scriptsize,
    breaklines=true,
    frame=single,
    showstringspaces=false
}

% --- Pie de página personalizado ---
\makeatletter

\setbeamertemplate{footline}{
  \leavevmode%
  \hbox{%
  \begin{beamercolorbox}[wd=.333333\paperwidth,ht=2.25ex,dp=1ex,center]{author in head/foot}%
    \usebeamerfont{author in head/foot}\insertshortauthor
  \end{beamercolorbox}%
  \begin{beamercolorbox}[wd=.333333\paperwidth,ht=2.25ex,dp=1ex,center]{title in head/foot}%
    \usebeamerfont{title in head/foot}Sesión 9
  \end{beamercolorbox}%
  \begin{beamercolorbox}[wd=.333333\paperwidth,ht=2.25ex,dp=1ex,right]{date in head/foot}%
    \usebeamerfont{date in head/foot}NLP \hfill \insertframenumber/\inserttotalframenumber \hspace*{0.3cm}
  \end{beamercolorbox}}%
  \vskip0pt%
}

\makeatother

% --- Datos de la portada ---
\title[SESIÓN 9 NLP]{\textbf{Sesión 9}}
\subtitle{Transformers y Atención Multi-Head}
\author[Prof. C. Sánchez]{Carlos César Sánchez Coronel}
\institute{\textbf{NLP}}
\date{}

% Logo opcional (descomentar si tienes la imagen)
% \titlegraphic{\includegraphics[height=1.5cm]{logo.png}}

% --- Eliminar la barra de navegación (opcional) ---
\setbeamertemplate{navigation symbols}{}

% --- Índice al inicio de cada sección ---
\AtBeginSection[]{
  \begin{frame}
    \frametitle{Contenido}
    \tableofcontents[currentsection]
  \end{frame}
}

% ============================================================================
% INICIO DEL DOCUMENTO
% ============================================================================
\begin{document}

% --- Portada ---
\begin{frame}
    \titlepage
\end{frame}

% --- Índice general ---
\begin{frame}{Contenido}
    \tableofcontents
\end{frame}

% ============================================================================
% SECCIÓN 1: INTRODUCCIÓN
% ============================================================================
\section{Introducción}

\begin{frame}{El fin de las recurrencias}
    \begin{itemize}
        \item Hasta 2017, las RNN/LSTM dominaban el NLP.
        \item Limitaciones:
        \begin{itemize}
            \item \textbf{Procesamiento secuencial}: no paralelizable, lentas en GPU.
            \item \textbf{Dependencias largas}: a pesar de LSTM, aún sufren gradiente desvaneciente.
            \item \textbf{Memoria limitada}: estado oculto comprime toda la información.
        \end{itemize}
        \item \textbf{Solución}: Transformers (Vaswani et al., 2017) – ``Attention is All You Need''.
    \end{itemize}
    % Basado en Pregunta 1
\end{frame}

\begin{frame}{Ventajas de los Transformers}
    \begin{itemize}
        \item \textbf{Paralelización total}: todas las posiciones se procesan simultáneamente.
        \item \textbf{Conexiones directas}: cada palabra atiende a cualquier otra en una sola capa.
        \item \textbf{Escalabilidad}: modelos con miles de millones de parámetros (GPT, BERT).
    \end{itemize}
\end{frame}

% ============================================================================
% SECCIÓN 2: ARQUITECTURA GENERAL
% ============================================================================
\section{Arquitectura General del Transformer}

\begin{frame}{Visión general}
    \begin{columns}[t]
        \column{0.5\textwidth}
        \textbf{Encoder}
        \begin{itemize}
            \item N bloques idénticos (típicamente 6).
            \item Cada bloque: Multi-Head Self-Attention + Feed-Forward Network.
            \item Conexiones residuales y layer norm.
        \end{itemize}
        \column{0.5\textwidth}
        \textbf{Decoder}
        \begin{itemize}
            \item N bloques idénticos.
            \item Añade una capa de Cross-Attention.
            \item Self-Attention enmascarada (para autoregresividad).
        \end{itemize}
    \end{columns}
    % Basado en Pregunta 2
    % IMAGEN: Diagrama del transformer
\end{frame}

% ============================================================================
% SECCIÓN 3: SELF-ATTENTION
% ============================================================================
\section{Self-Attention}

\begin{frame}{Intuición de self-attention}
    \begin{itemize}
        \item Cada palabra calcula la relevancia de las demás para actualizar su representación.
        \item Analogía: reunión donde cada persona (token) pregunta a las demás qué tan relevantes son.
    \end{itemize}
    % Basado en Pregunta 3
\end{frame}

\begin{frame}{Formulación matemática}
    Para una secuencia de $T$ palabras representada por $\mathbf{X} \in \mathbb{R}^{T \times d}$:
    \begin{enumerate}
        \item Proyectar en Queries, Keys, Values:
        \[
        \mathbf{Q} = \mathbf{X}\mathbf{W}^Q,\quad \mathbf{K} = \mathbf{X}\mathbf{W}^K,\quad \mathbf{V} = \mathbf{X}\mathbf{W}^V
        \]
        \item Calcular puntuaciones: $\mathbf{S} = \mathbf{Q}\mathbf{K}^T \in \mathbb{R}^{T \times T}$
        \item Escalar: $\mathbf{S}_{\text{scaled}} = \frac{\mathbf{S}}{\sqrt{d_k}}$
        \item Normalizar con softmax: $\mathbf{A} = \text{softmax}(\mathbf{S}_{\text{scaled}})$
        \item Salida: $\mathbf{Z} = \mathbf{A}\mathbf{V}$
    \end{enumerate}
    % Basado en Pregunta 3
\end{frame}

\begin{frame}{Escalado del producto punto}
    \begin{itemize}
        \item $\frac{1}{\sqrt{d_k}}$ evita que las puntuaciones sean demasiado grandes.
        \item Si los productos punto son grandes, el softmax se satura y los gradientes se vuelven muy pequeños.
        \item Mejora la estabilidad numérica.
    \end{itemize}
\end{frame}

% ============================================================================
% SECCIÓN 4: MULTI-HEAD ATTENTION
% ============================================================================
\section{Multi-Head Attention}

\begin{frame}{Motivación}
    \begin{itemize}
        \item Una sola cabeza puede promediar demasiado.
        \item Diferentes cabezas pueden aprender diferentes tipos de relaciones:
        \begin{itemize}
            \item Sintácticas (sujeto-verbo).
            \item Semánticas (sinónimos).
            \item De distancia.
        \end{itemize}
    \end{itemize}
    % Basado en Pregunta 4
\end{frame}

\begin{frame}{Formulación}
    \begin{enumerate}
        \item Para cada cabeza $i$: $\mathbf{Q}_i = \mathbf{X}\mathbf{W}_i^Q$, $\mathbf{K}_i = \mathbf{X}\mathbf{W}_i^K$, $\mathbf{V}_i = \mathbf{X}\mathbf{W}_i^V$.
        \item Calcular atención por cabeza: $\mathbf{Z}_i = \text{Attention}(\mathbf{Q}_i, \mathbf{K}_i, \mathbf{V}_i)$.
        \item Concatenar: $\mathbf{Z} = \text{Concat}(\mathbf{Z}_1, \dots, \mathbf{Z}_h)$.
        \item Proyectar: $\mathbf{Z}_{\text{final}} = \mathbf{Z}\mathbf{W}^O$.
    \end{enumerate}
    % Basado en Pregunta 4
\end{frame}

% ============================================================================
% SECCIÓN 5: POSITIONAL ENCODING
% ============================================================================
\section{Positional Encoding}

\begin{frame}{Necesidad de positional encoding}
    \begin{itemize}
        \item Self-attention es invariante a la posición: permutar las palabras produce las mismas representaciones.
        \item Para inyectar información de orden, se suman positional encodings a los embeddings de entrada.
    \end{itemize}
    % Basado en Pregunta 5
\end{frame}

\begin{frame}{Fórmula de positional encoding}
    \[
    \begin{align*}
        PE_{(pos, 2i)} &= \sin\left(\frac{pos}{10000^{2i/d}}\right) \\
        PE_{(pos, 2i+1)} &= \cos\left(\frac{pos}{10000^{2i/d}}\right)
    \end{align*}
    \]
    \begin{itemize}
        \item Funciones periódicas que permiten extrapolar a longitudes no vistas.
        \item No añaden parámetros entrenables.
    \end{itemize}
    % Basado en Pregunta 5
\end{frame}

% ============================================================================
% SECCIÓN 6: CAPA FEED-FORWARD
% ============================================================================
\section{Feed-Forward Network (FFN)}

\begin{frame}{FFN en transformers}
    \begin{block}{Arquitectura}
        \[
        \text{FFN}(x) = \max(0, x\mathbf{W}_1 + b_1)\mathbf{W}_2 + b_2
        \]
        \begin{itemize}
            \item Capa lineal con expansión (típicamente $4 \times d_{\text{model}}$).
            \item Activación ReLU (o GELU).
            \item Misma FFN para todas las posiciones, pero con parámetros compartidos.
        \end{itemize}
    \end{block}
    \begin{itemize}
        \item Aporta no linealidad y capacidad de transformación por token.
    \end{itemize}
    % Basado en Pregunta 10
\end{frame}

% ============================================================================
% SECCIÓN 7: CONEXIONES RESIDUALES Y LAYER NORM
% ============================================================================
\section{Conexiones Residuales y Layer Normalization}

\begin{frame}{Conexiones residuales}
    \begin{block}{Fórmula}
        \[
        x' = x + \text{Sublayer}(x)
        \]
    \end{block}
    \begin{itemize}
        \item Facilitan el flujo del gradiente en redes profundas.
        \item Permiten apilar muchas capas (ej. BERT Large: 24 capas).
    \end{itemize}
    % Basado en Pregunta 9
\end{frame}

\begin{frame}{Layer Normalization}
    \begin{block}{Fórmula}
        \[
        \text{LayerNorm}(x) = \frac{x - \mu}{\sigma} \odot \gamma + \beta
        \]
    \end{block}
    \begin{itemize}
        \item Normaliza las activaciones a lo largo de la dimensión de características.
        \item Independiente del batch size, funciona bien con secuencias de longitud variable.
    \end{itemize}
    % Basado en Pregunta 9
\end{frame}

% ============================================================================
% SECCIÓN 8: BLOQUES COMPLETOS
% ============================================================================
\section{Bloques Encoder y Decoder}

\begin{frame}{Encoder block}
    \begin{enumerate}
        \item Multi-Head Self-Attention con conexión residual + layer norm.
        \item Feed-Forward Network con conexión residual + layer norm.
    \end{enumerate}
\end{frame}

\begin{frame}{Decoder block}
    \begin{enumerate}
        \item \textbf{Masked Multi-Head Self-Attention}: evita que el decoder vea tokens futuros (máscara triangular).
        \item \textbf{Cross-Attention}: queries del decoder, keys/values del encoder.
        \item Feed-Forward Network.
    \end{enumerate}
    Cada subcapa con residual y layer norm.
    % Basado en Preguntas 7 y 8
\end{frame}

% ============================================================================
% SECCIÓN 9: COMPARACIÓN CON RNN
% ============================================================================
\section{Comparación Transformer vs RNN/LSTM}

\begin{frame}{Tabla comparativa}
    \begin{table}
        \centering
        \begin{tabular}{l|c|c}
            \toprule
            \textbf{Característica} & \textbf{Transformer} & \textbf{RNN/LSTM} \\
            \midrule
            Paralelización & Total (todas las posiciones) & Secuencial \\
            Dependencias largas & Conexiones directas & Gradiente desvaneciente \\
            Memoria & Acceso directo a todas las posiciones & Estado oculto comprimido \\
            Entrenamiento & Rápido en GPU & Lento \\
            Escalabilidad & Miles de millones de parámetros & Difícil de escalar \\
            \bottomrule
        \end{tabular}
    \end{table}
    % Basado en Pregunta 6
\end{frame}

% ============================================================================
% SECCIÓN 10: MODELOS DERIVADOS
% ============================================================================
\section{Modelos Derivados}

\begin{frame}{BERT (Bidirectional Encoder Representations from Transformers)}
    \begin{itemize}
        \item Usa solo el encoder (bidireccional).
        \item Entrenado con Masked Language Model (MLM) y Next Sentence Prediction (NSP).
        \item Ideal para tareas de comprensión (clasificación, NER, QA).
    \end{itemize}
    % Basado en Preguntas 11 y 13
\end{frame}

\begin{frame}{GPT (Generative Pre-trained Transformer)}
    \begin{itemize}
        \item Usa solo el decoder (autoregresivo).
        \item Entrenado con modelado de lenguaje (predecir la siguiente palabra).
        \item Ideal para generación de texto (ChatGPT).
    \end{itemize}
    % Basado en Pregunta 12
\end{frame}

\begin{frame}{T5 (Text-to-Text Transfer Transformer)}
    \begin{itemize}
        \item Arquitectura encoder-decoder.
        \item Formula todas las tareas como texto a texto.
        \item Ventaja: mismo pre-entrenamiento y fine-tuning para múltiples tareas.
    \end{itemize}
    % Basado en Pregunta 15
\end{frame}

\begin{frame}{Vision Transformer (ViT)}
    \begin{itemize}
        \item Aplica transformers a imágenes dividiéndolas en parches.
        \item Cada parche se proyecta linealmente y se añade positional encoding.
        \item Captura relaciones globales desde las primeras capas.
    \end{itemize}
    % Basado en Pregunta 16
\end{frame}

% ============================================================================
% SECCIÓN 11: IMPLEMENTACIÓN CONCEPTUAL
% ============================================================================
\section{Implementación Conceptual}

\begin{frame}[fragile]{Self-attention en Python}
    \begin{lstlisting}
import numpy as np

def softmax(x, axis=-1):
    e_x = np.exp(x - np.max(x, axis=axis, keepdims=True))
    return e_x / np.sum(e_x, axis=axis, keepdims=True)

def self_attention(X, W_q, W_k, W_v):
    # X: (T, d)
    Q = X @ W_q  # (T, d_k)
    K = X @ W_k  # (T, d_k)
    V = X @ W_v  # (T, d_k)
    scores = Q @ K.T  # (T, T)
    scaled = scores / np.sqrt(Q.shape[-1])
    weights = softmax(scaled)  # (T, T)
    output = weights @ V  # (T, d_k)
    return output, weights

# Ejemplo
T, d, d_k = 4, 8, 4
X = np.random.rand(T, d)
W_q = np.random.rand(d, d_k)
W_k = np.random.rand(d, d_k)
W_v = np.random.rand(d, d_k)
out, attn = self_attention(X, W_q, W_k, W_v)
print("Matriz de atención:\n", attn)
    \end{lstlisting}
\end{frame}

% ============================================================================
% SECCIÓN 12: PREPARACIÓN PARA ENTREVISTAS
% ============================================================================
\section{Preparación para Entrevistas}

\begin{frame}{Preguntas estrella}
    \begin{itemize}
        \item ¿Cómo funciona el transformer y qué problema resuelve?
        \item ¿Qué es positional encoding y por qué es necesario?
        \item Diferencia entre self-attention y cross-attention.
        \item Explica multi-head attention.
        \item ¿Por qué se escala el producto punto en atención?
    \end{itemize}
    % Basado en PDF
\end{frame}

\begin{frame}{Preguntas avanzadas}
    \begin{itemize}
        \item ¿Qué es la máscara en el decoder?
        \item ¿Cómo se calcula la complejidad computacional de la atención? $O(T^2 \cdot d)$
        \item ¿Qué ventajas tiene la atención relativa?
        \item ¿Cómo manejan contextos largos modelos como GPT-4? (atención dispersa, attention sinks)
        \item ¿Qué técnicas de optimización se usan para desplegar transformers en móviles? (pruning, quantization, distillation)
    \end{itemize}
    % Basado en Preguntas 7, 13, 14, 18, 19
\end{frame}

% ============================================================================
% SECCIÓN 13: CONEXIÓN CON EL RESTO DEL CURSO
% ========================================================

% ============================================================================
% SECCIÓN 14: RESUMEN
% ============================================================================
\section{Resumen}

\begin{frame}{Lo que hemos aprendido}
    \begin{exampleblock}{Conceptos clave}
        \begin{itemize}
            \item \textbf{Transformer}: arquitectura basada únicamente en atención, sin recurrencia.
            \item \textbf{Self-attention}: cada palabra atiende a todas las demás.
            \item \textbf{Multi-head attention}: varias cabezas capturan diferentes relaciones.
            \item \textbf{Positional encoding}: inyecta información de orden.
            \item \textbf{Encoder-decoder}: bloques con atención, FFN, residual y layer norm.
            \item \textbf{Ventajas}: paralelización, dependencias largas, escalabilidad.
            \item \textbf{Modelos derivados}: BERT (encoder), GPT (decoder), T5 (encoder-decoder), ViT (imágenes).
            \item \textbf{Desafíos}: complejidad cuadrática, inferencia en dispositivos pequeños.
        \end{itemize}
    \end{exampleblock}
\end{frame}

% --- Diapositiva final ---
\begin{frame}
    \centering
    \Huge \textbf{¡Gracias!}

\end{frame}

\end{document}